Une pile met en jeu les couples d'oxydoréduction
$\ce{Ni^{2+}_{(aq)}\ttslash{}Ni_{(s)}}$ et
$\ce{Cu^{2+}_{(aq)}\ttslash{}Cu_{(s)}}$, son schéma conventionnel est~:
$\circletlineh~\ce{Ni_{(s)}}~|~\ce{Ni^2+_{(aq)}}~\parallel~
\ce{Cu^2+_{(aq)}}~|~\ce{Cu_{(s)}}~\circletlinevh$.

Le nickel $\ce{Ni_{(s)}}$ est en excès et la quantité de matière initiale en
ions cuivre(II) est ${n_{i}(\ce{Cu^{2+}_{(aq)}})=\qty{1,0e-2}{\mole}}$.
L'intensité du courant débité par la pile est $I=\qty{40}{\mA}$ pendant toute la
durée de son fonctionnement.

\begin{donnees}
    \begin{itemize}
        \item $1~\mathcal{F}=\qty{9,65e4}{\coulomb\per\mole}$.
    \end{itemize}
\end{donnees}
\begin{enumerate}
    \item Écrire l'équation de la réaction de fonctionnement de la pile.
    \item Calculer $Q_{\max}$ la quantité d'électricité maximale débitée par la
          pile.
    \item Déterminer la durée de vie $\Delta t$ de cette pile.
\end{enumerate}

